\FloatBarrier
\section{Related work}
\label{sec:related}

\paragraph{Polar momentum and input-side preconditioning.}
\mname{Muon} introduced Newton--Schulz orthogonalization of matrix momentum~\citep{jordan2024muon}, motivating recent analyses of singular-value flattening, update norms, and subspace alignment~\citep{shumaylov2026freon,goyal2026spectrally}. Our counterfactual probes build on this separation between step norm and update direction. On the input side, \mname{Newton-Muon} studies when polar orthogonalization cancels left curvature~\citep{du2026newtonmuon}, \mname{MuonT} (Toner) introduces decoupled $(-1/2,-1/8)$ input activation inverse powers with cached-eigenbasis tracking~\citep{mohamadi2026toner}, and \mname{MuonH} replaces decoupled weight decay with Frobenius sphere retraction~\citep{wen2026hyperball} (\Cref{sec:hyperball}). \mname{Muon2-ExpDiag} extends \mname{MuonT} to make output-side preconditioning fast and measurable.

\paragraph{Kronecker, Gram, and two-sided polar preconditioners.}
Two-sided preconditioning has deep roots outside polar optimizers: \mname{K-FAC} factors curvature into activation and cotangent second moments~\citep{martens2015kfac}, while \mname{Shampoo} and \mname{SOAP} maintain left and right gradient Gram preconditioners~\citep{gupta2018shampoo,anil2020scalable,vyas2024soap}. Concurrent polar methods such as \mname{GO-MUON} and \mname{Mousse} apply symmetric moment or gradient-Gram inverse powers around orthogonalization with norm grafting~\citep{che2026gomuon,zhang2026mousse}. Rather than treating every projection alike, we assign dense or diagonal output factors by projection role, retain right-only weight decay, overlap output roots at a one-step lag, and test against execution-path controls, without claiming that cotangent moments outperform gradient statistics across all architectures or tuning regimes.

\section{Conclusion and scope}
\label{sec:conclusion}

\paragraph{Summary of contributions.}
\mname{Muon2-ExpDiag} shows that output-cotangent preconditioning improves a \mname{Muon} optimizer that already whitens its inputs, provided factor fidelity and scheduling match Transformer block geometry. Dense $(-1/2,-1/8)$ output roots at hidden width supply the cross-channel rotations needed by residual-stream writers, exact diagonal scaling on gated SwiGLU expansion outputs removes 96.2\% of nominal cubic output work where anisotropy is coordinate-aligned, right-only weight decay avoids late-cooldown degradation from output-shaped shrinkage, and a one-step side-stream lag overlaps output reductions and root repairs with training. Across three fresh 130M seeds, fourfold horizon scaling to 10.74B tokens, and width transfer to 520M, \mname{Muon2-ExpDiag} lowers held-out and out-of-distribution losses below tuned \mname{MuonT}, whereas all-diagonal output scaling (\mname{Muon2-Diag}) adds almost no held-out gain and regresses out of distribution at greater width.

\paragraph{Limitations and future directions.}
Several scope limits remain. The 80-TPP and 520M scaling runs each use one seed, all runs study a single pre-RMSNorm SwiGLU architecture family up to 537M non-embedding parameters, and our counterfactual and subspace-drift probes remain descriptive checkpoint audits rather than a causal proof of generalization. Periodic eigenbasis repairs also keep mean step-time overhead above median overhead ($10.7\text{--}17.0\%$ vs.\ $6.8\text{--}11.1\%$ over \mname{Muon}, or $8.9\text{--}11.4\%$ vs.\ $3.3\text{--}5.6\%$ above \mname{MuonT}), so fixed-token loss reductions do not yet yield wall-clock speedups. Head-block output approximations, adaptive root-refresh schedules, and multi-seed evaluations at multi-billion-parameter scale are natural next steps.
